\documentclass[18pt]{article}
\usepackage{amsmath,amsfonts,amssymb}
\usepackage{natbib}
\usepackage{booktabs}
\usepackage[utf8]{inputenc}
\usepackage[T1]{fontenc}
\usepackage{amssymb}
\usepackage[version=4]{mhchem}
\usepackage{stmaryrd}
\usepackage{bbold}
\usepackage{nopageno}
\usepackage{hyperref}
\usepackage{bookmark}
\author{Adam Handwerger}
\pagestyle{plain}

\begin{document}
\begin{flushleft}
Exercise 1
\vspace{.4cm}
Show that the Voronoi classifier can be written as\\
$f(x)=\frac{1}{2}\left(\parallel x-C_{-} \parallel^2-\parallel x-C_{+}\parallel^2\right)$\\
\vspace{.4cm}
Defining $C = (1/2) \left( C_{+} + C_{-} \right)$ and expanding this equation gives\\
\vspace{.4cm}
$f=\frac{1}{2} \left( \langle x-C_{-},x-C_{-}\rangle - \langle x-C_{+},x-C_{+}\rangle\right)=$
    $\frac{1}{2}\left( \langle x,x \rangle^2 -2 \langle x,C_{-} \rangle + \langle C_{-},C_{-}\rangle-\langle x,x \rangle^2 +2\langle x,C_{+} \rangle-\langle C_{+},C_{+}\rangle\right)=$\\
    $\frac{1}{2}\left(2\langle x,C_{+}\rangle-2\langle x,C_{-}\rangle+\langle C_{-},C_{-}\rangle-\langle C_{+},C_{+}\rangle\right)=$\\
    $\langle x,C_{+} \rangle - \langle x,C_{-}\rangle+ \frac{1}{2}\langle C_{-},C_{-}\rangle-\frac{1}{2}\langle C_{+},C_{+}\rangle=$\\
    $\langle x, C_{+}-C_{-} \rangle +\frac{1}{2}\langle C_{-},C_{-}\rangle-\frac{1}{2}\langle C_{+},C_{+}\rangle=$\\
    $\langle x-C+C,C_{+}-C_{-}\rangle+\frac{1}{2}\langle C_{-},C_{-}\rangle-\frac{1}{2}\langle C_{+},C_{+}\rangle=$\\
    $\langle x-C,C_{+}-C_{-}\rangle+ \langle C,C_{+}-C_{-}\rangle+\frac{1}{2}\langle C_{-},C_{-}\rangle-\frac{1}{2}\langle C_{+},C_{+}\rangle$\\
\vspace{.4cm}
Now the last two terms cancell since\\
$\langle C,C_{+}-C_{-}\rangle+\frac{1}{2}\langle C_{-},C_{-}\rangle-\frac{1}{2}\langle C_{+},C_{+}\rangle=$\\
$\frac{1}{2} \langle C_{+}, C_{+}-C_{-} \rangle + \frac{1}{2} \langle C_{-}, C_{+}-C_{-} \rangle+\frac{1}{2}\langle C_{-},C_{-}\rangle-\frac{1}{2} \langle C_{+},C_{+}\rangle=$\\
$\frac{1}{2}\langle C_{+}, -C_{-}\rangle + \frac{1}{2} \langle C_{+}, C_{-}\rangle=0$\\
\vspace{.4cm}
Excercise2
\vspace{.4cm}
For the Voronoi classifier example, the kernelized f is given by\\
$f(x)=\sum_i \alpha_i k(x,x_i)$, with


\begin{align}
    \alpha_i=  \left\{
    \begin{aligned}
                &-(y_i/n_n),\text{if $x_i$ was misclassified as postive in training}\\
                & (y_i/n_p),\text{if $x_i$ was misclassified as negative in training}\\
                & 0, otherwise\\
    \end{aligned}
    \right.
\end{align}
\hspace{1.2cm} $b=\frac{1}{2}\left(\frac{1}{n^2_n} \sum_{i|y_i<0} \sum_{j|y_j<0} k(x_i,x_j)-\frac{1}{n^2_p}\sum_{i|y_i>0}\sum_{j|y_j>0} k(x_i,x_j)\right)$\\
\vspace{.4cm}
Exercise 3 \hspace{.5cm}   Show that $f_\phi(x) =\frac{1}{2}\left(\parallel \phi_x-\phi_{n}\parallel_H-\parallel \phi_x-\phi_{p}\parallel_H\right)$\\
\vspace{.4cm}
Following the details from Exercise 1\\
\vspace{.4cm}
If we take $\phi_{x*}=\phi_{C+-C-}$ and $b=\frac{1}{2}\left(K(C_{+},C_{+})^2-K \left(C_{-},C_{-}\right)^2\right)$, then
the sums over $x_i$ are taken before the mapping $\phi:\chi \to H$ and f takes the form\\
\vspace{.4cm}
$\hspace{4cm} f_\phi(x)=\langle \phi_x, \phi_{x*} \rangle+b$.\\
\vspace{.4cm}
A property that $\phi(x)$ should have is $\phi(x):x\rightarrow h(x-C)$ so that h does shifting\\
by C in calculating the inner product. In this way we model the defintion of f from exercise 1.\\
\pagebreak
Exercise 4

$\chi$ = $\{(e2,e2),(e2,e3),(e3,e3),(e3,e4),(e4,e4)\}$.\\

Set each element of $\chi$ to be the vector $x_i$, for $i \in \{1,\dots,5\}$.\\
\end{flushleft}
The kernel matrix, K given by
$$
K =
     {\begin{bmatrix}
    2&1&1&0&0 \\
    1&2&1&1&0 \\
    0&1&1&2&1 \\
    0&0&0&1&2 
    \end{bmatrix}}
$$
\vspace{.5cm}
\begin{flushleft}
Its eigenvalues are $\lambda $=(4.3027756,2.6180340,2.0000000,0.6972244,0.3819660).\\
\vspace{.5cm}
Prove that K is a positive definite kernel.\\
\vspace{.5cm}
K is positive definite since all its eigenvalues are positive.\\
\vspace{.5cm}
Exercise 5\\
Find a Hilbert Space H and a mapping $\phi$ such that $K(x,y)=\langle \phi_x,\phi_y \rangle $.
Consider $H=\mathbb{R}^6$ and set
\end{flushleft}

$$
\begin{aligned}
&\phi(x_1) = \left(1, 1, 0, 0, 0, 0\right) \\
&\phi(x_2) = \left(0, 1, 1, 0, 0, 0\right) \\
&\phi(x_3) = \left(0, 0, 1, 1, 0, 0\right) \\
&\phi(x_4) = \left(0, 0, 0, 1, 1, 0\right) \\
&\phi(x_5) = \left(0, 0, 0, 0, 1, 1\right) \\
\end{aligned}
$$

\begin{flushleft}
Then for any $x \text{ and } y \in \chi$, every element of K can be represented as\\
$k(x,y)=\langle \phi(x),\phi(y)\rangle$.\\
\vspace{.5cm}
Exercise 6
\\
Show that $k(x,y)=x^{\top}y$ is pos. definite.\\

$$
\begin{aligned}
  &v^{\top} K v=\sum_{i j}^n v_{i} v_{j} k_(x_i,x_j)=\sum_{i j} v_{i} v_{j}\left(x_{i}^{\top} x_{j}\right) \\
= & \sum_{i} \sum_{j} v_{i} v_{j} x_{i}^{\top} x_{j} = \left(\sum_{i} v_{i} x_{i}^{\top}\right)\left(\sum_{j} v_{j} x_{j}\right) \\
= & \left(\sum_{i} v_{i} x_{i}\right)^{\top}\left(\sum_{j} v_{j} x_{j}\right)  \geqslant 0\\
\end{aligned}
$$
\vspace{.4cm}
If set $w=\sum v_i x_i$ then last term is $\Vert w\Vert$  which is greater or equall to zero.
Assume $x\in\mathbb{R}^d$ and set $k(x,y)=x^{\top}y$.\\
\vspace{.4cm}
The construction of the associated Hilbert space will follow discussion in class.\\
\vspace{.4cm}
$f(x)=\sum_i^n k(x,x_i)=\sum_i \alpha_i x^{\top}x_i=\sum_i\alpha_ix_i^{\top}x$=\\
\vspace{.4cm}
$\left(\sum_i\alpha_i x_i\right)^{\top}x=v^{\top}x$\\
\vspace{.4cm}
Guess: H = $\{f_v | f_v(x)=v^{\top}x, v \in \mathbb{R}^d\}$.\\
\vspace{.4cm}
Check to see if this construction has the correct properties:\\
\vspace{.4cm}
0. Since by construction $v \in\mathbb{R}^d$, the Hilbert Space is defined\\
on vectors of finite dimension and is hence is complete.\\
\vspace{.4cm}
1. Now we check to see if $k(x,y)=x^{\top}y$ implies an f of the form $f_v$.\\
\vspace{.4cm}
Let $x\in\chi$. Since $x\longrightarrow x^{\top}y=y^{\top}x=f(x)$ for $y\in \mathbb{R}^d,$\\
$k(x,y)$ will produce $f$ of the form $f_v\in H$.\\
\vspace{.4cm}
2. Check that the k is a reproducing kernel. I.e $ \langle f_v,k(\cdot,y)\rangle=f_v(y)$
\vspace{.4cm}
Since $f_v(y)=v^{\top}y$,\\
$\langle f_v,k(\cdot ,y)\rangle =\langle f_v,f_y\rangle=v^{\top}y=f_v(y)$\\
\vspace{.4cm}
So k is reproducing.\\
\vspace{.5cm}
Excersize 7\\
\vspace{.4cm}
The quadratic kernel given by $\left(x^Ty\right)^2$ is pos. definite.\\
\vspace{.4cm}
Proof: Take any vector $v\in\mathbb{R}^n$. Then the quadratic form is given by\\

$v^{\top}Kv=\sum_{ij}\alpha_i \beta_j K_{ij}=\sum_{ij} \alpha_i \beta_j \langle \phi(x_i),\phi(x_j)\rangle_H$.
Since for any pair of points with components given by $x=(x_1,x_2)$ and $x^{\prime}=(x^{\prime}_1,x^{\prime}_2)$,\\ the kernel reduces to $\left(x_1 x_1^{\prime}+_2 x_2^{\prime}\right)^2$ and each element is symtetric
with respect to the primed coordinates. This means K is as well. Then by construction K can be decomposed as $K=A^{\top}A$ and the quadratic form becomes\\
\vspace{.4cm}
$v^{\top}Kv=v^{\top}A A^{\top}v=\left(v^{\top}A\right)\left(A^{\top}v \right)=\left(A^{\top}v\right)^{\top} A^{\top}v=\parallel Av \parallel^2 \geq 0$\\




\end{flushleft}
\end{document}


